\RequirePackage{iftex}
\ifPDFTeX
\documentclass[a4paper,12pt]{article}
\usepackage[margin=24mm]{geometry}
\usepackage[T1]{fontenc}
\usepackage[utf8]{inputenc}
\usepackage{CJKutf8}
\else
\documentclass[a4paper,12pt]{jsarticle}
\usepackage[dvipdfmx,margin=24mm]{geometry}
\fi
\usepackage{amsmath,amssymb}
\usepackage{enumitem}

\setlength{\parindent}{0pt}
\setlength{\parskip}{0.35em}
\setlist[enumerate]{label=(\alph*),leftmargin=3em,itemsep=0.25em,topsep=0.35em}

\begin{document}
\ifPDFTeX
\begin{CJK}{UTF8}{min}
\fi

\newcommand{\problem}[2]{\par\vspace{1.0em}\noindent{\large 問題 #1}\quad（#2点）\quad\ignorespaces}
\newcommand{\choiceproblem}[2]{\par\vspace{1.0em}\noindent{\large 選択問題 #1}\quad（#2点）\quad\ignorespaces}

\begin{center}
{\Large 数学1A 中間試験}\\[0.4em]
\end{center}

計算問題では途中式を、証明問題ではどの定理を使ったか書くこと。各10点である。

\problem{1}{$5\times 2 = 10$}
次の文を，$\mathbb{R},\mathbb{N},\forall,\exists$ などの記号を用いて書き直せ。

\begin{enumerate}
\item 任意の実数 $x$ に対して，$x^2+1>0$ である。
\item ある自然数 $n$ が存在して，$n^2>100$ となる。
\end{enumerate}

\problem{2}{$5\times 2 = 10$}
(1)  次の関数が $x=1$ で連続となるような実数 $a$ を求めよ。
\[
f(x)=
\begin{cases}
x^2+a & (x<1),\\
2x & (x\ge 1).
\end{cases}
\]

(2) 関数 $f(x)=x^3+x-1$ を考える。中間値の定理を用いて，方程式 $x^3+x-1=0$ が区間 $(0,1)$ に少なくとも1つの解をもつことを示せ。

\problem{3}{$5\times 2 = 10$}
次の関数の微分を求めよ: (1) $f(x)=(x^2+1)\sin x$,
(2) $g(x)= x\log x $, $(x>0)$.

\problem{4}{10}
ロピタルの定理を用いて，次の極限を求めよ。
 $$\lim_{x\to 0}\dfrac{\cos (\sin{x})-1}{x}.$$

\problem{5}{10}
二変数関数 $f(x,y)=x^3y^2+e^{xy}$ について, 偏微分
$\frac{\partial f}{\partial x}(x,y),\,\frac{\partial f}{\partial y}(x,y)$を求めよ.


\problem{6}{10}
$f(x)=\sin x$ の $a=\dfrac{\pi}{2}$ における3次のテイラー多項式を求めよ。
\problem{7}{10}
 $f(x)=e^{2x}$ の マクローリン展開（つまり$a=0$のテイラー展開）を求めよ。

\problem{8}{$5\times 2=10$}
 指数関数 $e^x$ の $a=0$ における$n$次のテイラー多項式を $T_n$とおく。
次の手順に従って, $e$が無理数であることを示せ: 
\begin{enumerate}
\item テイラーの定理を用いて, $n>3$のとき次を示せ ($e<3$は用いてよい):
$$0<|e - T_n(1)|<\frac{1}{n!} .$$
\item (a)を用いて, $e$ が有理数であると仮定して矛盾を導け。
\end{enumerate}

\problem{9}{$5\times 2 = 10$}
\begin{enumerate}
\item オイラーの公式を用いて,　実数 $x$について $\cos{(i x)}$ を$e^x,e^{-x}$を用いて表せ。ここで, $i$ は虚数単位。
\item $\cos{z}=2$となる複素数$z$を一つ求めよ。
\end{enumerate}
\problem{10}{$5\times 2=10$}
\begin{enumerate}
    \item 関数$f(x) = 1+ \sin\Big(\frac{\pi}4 x\Big)$について, $f(x)$が$x\in [0,2]$で単調増加であることを示せ. (ただし$(0,2)$で$f$の微分が正ならば$[0,2]$上で単調増加であることは用いてよい).  
\item $a_0=0,$ $a_{n+1}=1+\sin{\left(\frac{\pi}{4} a_n \right)}$で定まる数列$(a_n)$の極限を求め、その値がなぜ極限になるか説明せよ。
\end{enumerate}
\ifPDFTeX
\end{CJK}
\fi

\end{document}
